% GENERATED FILE. Edit canonical lessons, then run npm run generate:books.
% canonical-source-hash: fnv1a64:1c10a85b440609b7
% canonical-lessons: PA-C294-kallar, PA-C294-rakkar, PA-C294-ujar, PA-C294-dhela, PA-C294-dhur, PA-R294-review-288-294

\chapter{Poor Ground}
\label{ch:pa-a2-poor-ground}
\hlchaptermodality{294}

\begin{chapteropening}
\textbf{By the end of this chapter:} \emph{I can say what kind of ground a place has.}

Complete the last lesson of chapter 294: I can say what kind of ground a place has.
\end{chapteropening}

\section[\texorpdfstring{\pa{ਕੱਲਰ} (kallar)}{ਕੱਲਰ (kallar)}]{\pa{ਕੱਲਰ} (kallar) — salty barren soil}

\label{lesson:PA-C294-kallar}



\begin{quote}
Before the new one: say the Punjabi for a slope, then the Punjabi for a hillock.
\end{quote}

\subsection*{You'll want to know: \pa{ਕੱਲਰ}}
\textbf{\pa{ਕੱਲਰ}} — \emph{kallar} — \textquotedblleft{}salty barren soil\textquotedblright{}.

\begin{quote}
\pa{ਕੱਲਰ} \pa{ਵਿੱਚ} \pa{ਕਣਕ} \pa{ਚੰਗੀ} \pa{ਨਹੀਂ} \pa{ਹੁੰਦੀ} \pa{ਹੈ।}
\end{quote}

\begin{quote}
\emph{kallar vich kaṇak chaṅgī nahīṁ hundī hai.}
\end{quote}

\begin{quote}
Wheat does not grow well in salty soil.
\end{quote}

A masculine noun.

\begin{cousinweb}
\textbf{\pa{ਕੱਲਰ}} (\emph{kallar}) is inherited from a reconstructed \emph{*kallara}, \textquotedblleft{}salt barren soil\textquotedblright{}; CDIAL lists Punjabi \emph{kallar}, \textquotedblleft{}salt barren soil\textquotedblright{}, under it (CDIAL 2954).
\end{cousinweb}

\begin{grammarlens}[title={What you've built}]
One more word for this chapter.
\end{grammarlens}

\subsection*{Your turn}
\begin{itemize}
\raggedright
  \item \emph{Say it:} \emph{kallar}
  \item \emph{Say it:} \emph{kallar}, once more
  \item \emph{Say it:} \emph{kallar vich kaṇak chaṅgī nahīṁ hundī hai.} — Wheat does not grow well in salty soil.
  \item \emph{Recall:} say the Punjabi for a slope, then the Punjabi for a hillock, then say \emph{kallar} again
\end{itemize}

\begin{tcolorbox}[breakable,colback=teal!4,colframe=teal!35!black,title={Before you move on}]
What does \pa{ਕੱਲਰ} mean? (\textquotedblleft{}Salty barren soil\textquotedblright{}.) Say it once more.
\end{tcolorbox}
\section[\texorpdfstring{\pa{ਰੱਕੜ} (rakkaṛ)}{ਰੱਕੜ (rakkaṛ)}]{\pa{ਰੱਕੜ} (rakkaṛ) — hard and barren (of land)}

\label{lesson:PA-C294-rakkar}



\begin{quote}
Before the new one: say the Punjabi for a hillock, then the Punjabi for salty barren soil.
\end{quote}

\subsection*{You'll want to know: \pa{ਰੱਕੜ}}
\textbf{\pa{ਰੱਕੜ}} — \emph{rakkaṛ} — \textquotedblleft{}hard and barren (of land)\textquotedblright{}.

\begin{quote}
\pa{ਇਹ} \pa{ਜ਼ਮੀਨ} \pa{ਰੱਕੜ} \pa{ਹੈ।}
\end{quote}

\begin{quote}
\emph{eh zamīn rakkaṛ hai.}
\end{quote}

\begin{quote}
This land is hard and barren.
\end{quote}

\begin{cousinweb}
\textbf{\pa{ਰੱਕੜ}} (\emph{rakkaṛ}) is inherited from a reconstructed \emph{*rakka}, \textquotedblleft{}defective\textquotedblright{}; CDIAL lists Punjabi \emph{rakkaṛ}, \textquotedblleft{}hard and barren (of land)\textquotedblright{}, under it (CDIAL 10538).
\end{cousinweb}

\begin{grammarlens}[title={What you've built}]
One more word for this chapter.
\end{grammarlens}

\subsection*{Your turn}
\begin{itemize}
\raggedright
  \item \emph{Say it:} \emph{rakkaṛ}
  \item \emph{Say it:} \emph{rakkaṛ}, once more
  \item \emph{Say it:} \emph{eh zamīn rakkaṛ hai.} — This land is hard and barren.
  \item \emph{Recall:} say the Punjabi for a hillock, then the Punjabi for salty barren soil, then say \emph{rakkaṛ} again
  \item \emph{Write it:} copy \textbf{\pa{ਰੱਕੜ}} once, checking the whole word
\end{itemize}

\begin{tcolorbox}[breakable,colback=teal!4,colframe=teal!35!black,title={Before you move on}]
What does \pa{ਰੱਕੜ} mean? (\textquotedblleft{}Hard and barren (of land)\textquotedblright{}.) Say it once more.
\end{tcolorbox}
\section[\texorpdfstring{\pa{ਉਜਾੜ} (ujāṛ)}{ਉਜਾੜ (ujāṛ)}]{\pa{ਉਜਾੜ} (ujāṛ) — desolate; deserted}

\label{lesson:PA-C294-ujar}



\begin{quote}
Before the new one: say the Punjabi for salty barren soil, then the Punjabi for barren.
\end{quote}

\subsection*{You'll want to know: \pa{ਉਜਾੜ}}
\textbf{\pa{ਉਜਾੜ}} — \emph{ujāṛ} — \textquotedblleft{}desolate; deserted\textquotedblright{}.

\begin{quote}
\pa{ਰਾਤ} \pa{ਨੂੰ} \pa{ਇਹ} \pa{ਸੜਕ} \pa{ਉਜਾੜ} \pa{ਹੁੰਦੀ} \pa{ਹੈ।}
\end{quote}

\begin{quote}
\emph{rāt nū̃ eh saṛak ujāṛ hundī hai.}
\end{quote}

\begin{quote}
At night this road is deserted.
\end{quote}

\begin{cousinweb}
\textbf{\pa{ਉਜਾੜ}} (\emph{ujāṛ}) is inherited from a reconstructed \emph{*ujjaṭa}, \textquotedblleft{}uprooted, laid waste\textquotedblright{}; CDIAL lists Punjabi \emph{ujāṛ}, \textquotedblleft{}waste, desolate\textquotedblright{}, under it (CDIAL 1660b).
\end{cousinweb}

\begin{grammarlens}[title={What you've built}]
One more word for this chapter.
\end{grammarlens}

\subsection*{Your turn}
\begin{itemize}
\raggedright
  \item \emph{Say it:} \emph{ujāṛ}
  \item \emph{Say it:} \emph{ujāṛ}, once more
  \item \emph{Say it:} \emph{rāt nū̃ eh saṛak ujāṛ hundī hai.} — At night this road is deserted.
  \item \emph{Recall:} say the Punjabi for salty barren soil, then the Punjabi for barren, then say \emph{ujāṛ} again
  \item \emph{Write it:} copy \textbf{\pa{ਉਜਾੜ}} once, checking the whole word
\end{itemize}

\begin{tcolorbox}[breakable,colback=teal!4,colframe=teal!35!black,title={Before you move on}]
What does \pa{ਉਜਾੜ} mean? (\textquotedblleft{}Desolate; deserted\textquotedblright{}.) Say it once more.
\end{tcolorbox}
\section[\texorpdfstring{\pa{ਢੇਲਾ} (ḍhelā)}{ਢੇਲਾ (ḍhelā)}]{\pa{ਢੇਲਾ} (ḍhelā) — a clod of earth}

\label{lesson:PA-C294-dhela}



\begin{quote}
Before the new one: say the Punjabi for barren, then the Punjabi for desolate.
\end{quote}

\subsection*{You'll want to know: \pa{ਢੇਲਾ}}
\textbf{\pa{ਢੇਲਾ}} — \emph{ḍhelā} — \textquotedblleft{}a clod of earth\textquotedblright{}.

\begin{quote}
\pa{ਸਾਡੇ} \pa{ਖੇਤ} \pa{ਵਿੱਚ} \pa{ਵੱਡੇ} \pa{ਢੇਲੇ} \pa{ਹਨ।}
\end{quote}

\begin{quote}
\emph{sāḍe khet vich vaḍḍe ḍhele han.}
\end{quote}

\begin{quote}
There are big clods in our field.
\end{quote}

A masculine noun.

\begin{cousinweb}
\textbf{\pa{ਢੇਲਾ}} (\emph{ḍhelā}) is inherited from a reconstructed \emph{*ḍala}, \textquotedblleft{}lump\textquotedblright{}, in its form \emph{*ḍhēlla}; CDIAL lists Punjabi \emph{ḍhelā}, \textquotedblleft{}clod\textquotedblright{}, under it (CDIAL 5536).
\end{cousinweb}

\begin{grammarlens}[title={What you've built}]
One more word for this chapter.
\end{grammarlens}

\subsection*{Your turn}
\begin{itemize}
\raggedright
  \item \emph{Say it:} \emph{ḍhelā}
  \item \emph{Say it:} \emph{ḍhelā}, once more
  \item \emph{Say it:} \emph{sāḍe khet vich vaḍḍe ḍhele han.} — There are big clods in our field.
  \item \emph{Recall:} say the Punjabi for barren, then the Punjabi for desolate, then say \emph{ḍhelā} again
  \item \emph{Write it:} copy \textbf{\pa{ਢੇਲਾ}} once, checking the whole word
\end{itemize}

\begin{tcolorbox}[breakable,colback=teal!4,colframe=teal!35!black,title={Before you move on}]
What does \pa{ਢੇਲਾ} mean? (\textquotedblleft{}A clod of earth\textquotedblright{}.) Say it once more.
\end{tcolorbox}
\section[\texorpdfstring{\pa{ਧੂੜ} (dhūṛ)}{ਧੂੜ (dhūṛ)}]{\pa{ਧੂੜ} (dhūṛ) — dust}

\label{lesson:PA-C294-dhur}



\begin{quote}
Before the new one: say the Punjabi for desolate, then the Punjabi for a clod.
\end{quote}

\subsection*{You'll want to know: \pa{ਧੂੜ}}
\textbf{\pa{ਧੂੜ}} — \emph{dhūṛ} — \textquotedblleft{}dust\textquotedblright{}.

\begin{quote}
\pa{ਸੜਕ} \pa{ਉੱਤੇ} \pa{ਬਹੁਤ} \pa{ਧੂੜ} \pa{ਹੈ।}
\end{quote}

\begin{quote}
\emph{saṛak utte bahut dhūṛ hai.}
\end{quote}

\begin{quote}
There is a lot of dust on the road.
\end{quote}

A feminine noun.

\begin{cousinweb}
\textbf{\pa{ਧੂੜ}} (\emph{dhūṛ}) is inherited from a reconstructed \emph{*dhūḍi}, \textquotedblleft{}dust, powder\textquotedblright{}; CDIAL lists Punjabi \emph{dhūṛ}, \textquotedblleft{}dust\textquotedblright{}, under it (CDIAL 6835).
\end{cousinweb}

\begin{grammarlens}[title={What you've built}]
One more word for this chapter.
\end{grammarlens}

\subsection*{Your turn}
\begin{itemize}
\raggedright
  \item \emph{Say it:} \emph{dhūṛ}
  \item \emph{Say it:} \emph{dhūṛ}, once more
  \item \emph{Say it:} \emph{saṛak utte bahut dhūṛ hai.} — There is a lot of dust on the road.
  \item \emph{Recall:} say the Punjabi for desolate, then the Punjabi for a clod, then say \emph{dhūṛ} again
  \item \emph{Write it:} copy \textbf{\pa{ਧੂੜ}} once, checking the whole word
\end{itemize}

\begin{tcolorbox}[breakable,colback=teal!4,colframe=teal!35!black,title={Before you move on}]
What does \pa{ਧੂੜ} mean? (\textquotedblleft{}Dust\textquotedblright{}.) Say it once more.
\end{tcolorbox}
\section[(dialogue)]{Looking back: from The Hot Season Ahead to Poor Ground}

\label{lesson:PA-R294-review-288-294}



\begin{quote}
Say the Punjabi for sunshine and a hot wind.
\end{quote}

\subsection*{Your turn}
Say these aloud:

\begin{itemize}
\raggedright
  \item sunshine, a hot wind, to blaze with heat, the hot season, a drought
  \item fog, frost, bitter cold, to be chilled, the cold season
  \item the sky, hail, to clear up, to go down, to blow
  \item cool, pleasantly warm, hot to the touch, lukewarm, to cool something down
  \item damp, arid, moisture, to dry up, to seep
  \item soil, land, level, a slope, a hillock
  \item salty barren soil, barren, desolate, a clod, dust
\end{itemize}

\begin{tcolorbox}[breakable,colback=teal!4,colframe=teal!35!black,title={Before you move on}]
Sunshine? (\textbf{\pa{ਧੁੱਪ}}, \emph{dhupp}.) A hot wind? (\textbf{\pa{ਲੂ}}, \emph{lū}.) To blaze with heat? (\textbf{\pa{ਤਪਣਾ}}, \emph{tapṇā}.) The hot season? (\textbf{\pa{ਹੁਨਾਲ}}, \emph{hunāl}.) A drought? (\textbf{\pa{ਸੋਕਾ}}, \emph{sokā}.) Fog? (\textbf{\pa{ਧੁੰਦ}}, \emph{dhund}.) Frost? (\textbf{\pa{ਕੱਕਰ}}, \emph{kakkar}.) Bitter cold? (\textbf{\pa{ਪਾਲਾ}}, \emph{pālā}.) To be chilled? (\textbf{\pa{ਠਰਨਾ}}, \emph{ṭharnā}.) The cold season? (\textbf{\pa{ਸਿਆਲ}}, \emph{siāl}.) The sky? (\textbf{\pa{ਅਸਮਾਨ}}, \emph{asmān}.) Hail? (\textbf{\pa{ਗੜੇ}}, \emph{gaṛe}.) To clear up? (\textbf{\pa{ਖੁੱਲ੍ਹਣਾ}}, \emph{khullhṇā}.) To go down? (\textbf{\pa{ਘਟਣਾ}}, \emph{ghaṭṇā}.) To blow? (\textbf{\pa{ਵਗਣਾ}}, \emph{vagṇā}.) Cool? (\textbf{\pa{ਠੰਢਾ}}, \emph{ṭhaṇḍhā}.) Pleasantly warm? (\textbf{\pa{ਨਿੱਘਾ}}, \emph{nigghā}.) Hot to the touch? (\textbf{\pa{ਤੱਤਾ}}, \emph{tattā}.) Lukewarm? (\textbf{\pa{ਕੋਸਾ}}, \emph{kosā}.) To cool something down? (\textbf{\pa{ਠਾਰਨਾ}}, \emph{ṭhārnā}.) Damp? (\textbf{\pa{ਸਿੱਲ੍ਹਾ}}, \emph{sillhā}.) Arid? (\textbf{\pa{ਖ਼ੁਸ਼ਕ}}, \emph{khushak}.) Moisture? (\textbf{\pa{ਨਮੀ}}, \emph{namī}.) To dry up? (\textbf{\pa{ਸੁੱਕਣਾ}}, \emph{sukkṇā}.) To seep? (\textbf{\pa{ਸਿੰਮਣਾ}}, \emph{simmṇā}.) Soil? (\textbf{\pa{ਮਿੱਟੀ}}, \emph{miṭṭī}.) Land? (\textbf{\pa{ਜ਼ਮੀਨ}}, \emph{zamīn}.) Level? (\textbf{\pa{ਪੱਧਰਾ}}, \emph{paddhrā}.) A slope? (\textbf{\pa{ਢਲਾਣ}}, \emph{ḍhalāṇ}.) A hillock? (\textbf{\pa{ਟਿੱਲਾ}}, \emph{ṭillā}.) Salty barren soil? (\textbf{\pa{ਕੱਲਰ}}, \emph{kallar}.) Barren? (\textbf{\pa{ਰੱਕੜ}}, \emph{rakkaṛ}.) Desolate? (\textbf{\pa{ਉਜਾੜ}}, \emph{ujāṛ}.) A clod? (\textbf{\pa{ਢੇਲਾ}}, \emph{ḍhelā}.) Dust? (\textbf{\pa{ਧੂੜ}}, \emph{dhūṛ}.)
\end{tcolorbox}
